\documentclass[pdflatex,sn-mathphys-num]{sn-jnl}

\usepackage{graphicx}
\usepackage{booktabs}
\usepackage{amsmath}
\usepackage{amssymb}
\usepackage{multirow}
\usepackage{longtable}

\begin{document}

\title{Joint Optimization of Hash-Based Signatures for Resource-Constrained Devices: Signature Compression, Compact State Management, and Performance Trade-offs}

\author[1]{\fnm{Sohaib} \sur{Rana}}
\email{sohaibrana2023@student.usm.my}

\author*[1,2]{\fnm{Mohd Najwadi} \sur{Yusoff}}
\email{najwadi@usm.my}

\author*[3,4]{\fnm{Je Sen} \sur{Teh}}
\email{j.teh@deakin.edu.au}

\author[3]{\fnm{Adnan} \sur{Anwar}}
\email{adnan.anwar@deakin.edu.au}

\affil[1]{%
\orgdiv{School of Computer Sciences},
\orgname{Universiti Sains Malaysia},
\orgaddress{\city{Gelugor}, \state{Penang}, \country{Malaysia}}}

\affil[2]{%
\orgdiv{Cybersecurity Research Centre (CYRES)},
\orgname{Universiti Sains Malaysia},
\orgaddress{\city{Gelugor}, \state{Penang}, \country{Malaysia}}}

\affil[3]{%
\orgdiv{School of Information Technology},
\orgname{Deakin University},
\orgaddress{\country{Australia}}}

% TODO: verify whether a fourth affiliation is required for Je Sen Teh.
% Do not guess the missing affiliation here.

\abstract{
Hash-based signatures provide post-quantum security with comparatively simple assumptions, but deployment on resource-constrained platforms remains sensitive to signature size, signing latency, state-management overhead, memory footprint, and implementation complexity. This study investigates a coordinated optimization framework for stateful hash-based signatures on resource-constrained systems. Individual and combined optimization strategies are evaluated using correctness-preserving implementation gates, controlled Raspberry Pi~5 measurements, and multi-objective analysis. The manuscript integrates experimental results across parameter optimization, compact state handling, alternative hash backends, memory-oriented optimization, parallel execution, signature-size reduction, cross-construction extensions, and combined optimization strategies. Final quantitative conclusions are inserted only from frozen validated datasets.
}

\keywords{
Hash-based signatures,
XMSS,
XMSSMT,
WOTS+,
post-quantum cryptography,
resource-constrained devices,
Raspberry Pi 5,
signature compression,
state management,
multi-objective optimization
}

\maketitle

\section{Introduction}
\label{sec:introduction}

% TODO:
% - Motivation for post-quantum signatures on constrained systems.
% - Signature-size and state-management problems.
% - Gap: individual optimizations are often studied separately.
% - Contribution: coordinated multi-objective evaluation.
% - Explicitly distinguish standardized constructions from
%   experimental non-standard extensions.

\section{Background}
\label{sec:background}

\subsection{XMSS and XMSSMT}
% TODO.

\subsection{WOTS+}
% TODO.

\subsection{State management in stateful hash-based signatures}
% TODO.

\subsection{Performance constraints on resource-constrained devices}
% TODO.

\section{Research Problem and Optimization Framework}
\label{sec:framework}

% TODO:
% Define the Objective-1 optimization dimensions:
% signature size,
% key generation latency,
% signing latency,
% verification latency,
% serialized state,
% memory,
% implementation cost.

\subsection{Optimization objectives}
% TODO.

\subsection{Experimental and standardized configurations}
% TODO:
% Clearly label every experimental/non-standard profile.

\section{Methodology}

% BEGIN PJ04_R2_METHODOLOGY
\subsection{PJ04 WOTS+C corrected R2 evaluation}

The PJ04 corrected R2 experiment compared the validated CONTROL and
PJ04\_WOTSC configurations under the measured Raspberry Pi~5 conditions.
Each mode comprised 10 execution blocks. Each block contained one key
generation operation, 10 signing operations, and 10 verification operations.
Accordingly, the 100 signing and 100 verification observations per mode were
nested within 10 execution blocks and are not interpreted as 100 independent
experimental replicates. CONTROL was executed before PJ04\_WOTSC; therefore,
temporal-order effects cannot be excluded.

\subsection{PJ04 block-aware uncertainty analysis}

Descriptive summaries were retained for the raw timing observations, whereas
uncertainty for signing and verification was evaluated at the execution-block
level. The validated nonparametric procedure independently resampled the 10
blocks within each mode, used 10,000 bootstrap resamples with seed 20261003,
and reported percentile 95\% confidence intervals. This independent-mode
resampling does not establish paired-block correspondence. The resulting
intervals are used as block-aware uncertainty summaries rather than as
confirmatory significance tests.
% END PJ04_R2_METHODOLOGY

\label{sec:methodology}

\subsection{Experimental platform}
% TODO:
% Raspberry Pi 5 hardware/software/environment.

\subsection{Measurement methodology}
% TODO:
% Timing clocks, affinity, thermal control, run-order control.

\subsection{Correctness and validation gates}
% TODO:
% Positive/negative testing, rollover, sanitizers, provenance.

\subsection{Statistical methodology}
% The M3-specific methodology is inserted below.
\input{m3_rpi5/M3_RPI5_METHODS_STATISTICAL_ANALYSIS.tex}

\section{Individual Optimization Strategies}

% BEGIN PJ04_R2_INDIVIDUAL_OPT
\subsection{PJ04 WOTS+C corrected R2 artifact and latency results}

PJ04\_WOTSC reduced the XMSS signature from 2500 to 2408~bytes, an exact
reduction of 92~bytes. The corresponding C1 FULL packet decreased from 2506
to 2414~bytes, also an exact 92-byte reduction. These byte differences are
deterministic artifact-geometry results and are therefore reported without
inferential significance testing.

\input{pj04_r2_tables/PJ04_R2_TABLE_ARTIFACT_GEOMETRY.tex}

\begin{figure*}[htbp]
\centering
\includegraphics[width=0.88\textwidth]{pj04_r2_assets/PJ04_R2_FIG_ARTIFACT_SIZE_REDUCTION.pdf}
\caption{PJ04 R2 deterministic signature and FULL-packet size comparison.}
\label{fig:pj04_r2_size}
\end{figure*}

Under the measured Raspberry Pi~5 conditions, the observed mean key-generation
latency changed by $-4.935397\%$ for PJ04\_WOTSC relative to CONTROL. The
block-aware bootstrap interval for the percentage change was
$[-5.169274124\%, -4.739936765\%]$.

The observed mean signing latency changed by $-4.833727\%$, with a block-aware
bootstrap interval of $[-4.934408761\%, -4.721328801\%]$. The observed mean
verification latency changed by $-3.800672\%$, with a corresponding interval
of $[-4.944207485\%, -2.468006971\%]$. These measurements describe the
observed implementation behavior under the evaluated conditions and are not
interpreted as causal performance effects or universal superiority.

\input{pj04_r2_tables/PJ04_R2_TABLE_LATENCY_SUMMARY.tex}

\input{pj04_r2_tables/PJ04_R2_TABLE_EFFECTS_AND_CI.tex}

\begin{figure*}[htbp]
\centering
\includegraphics[width=0.88\textwidth]{pj04_r2_assets/PJ04_R2_FIG_LATENCY_EFFECTS.pdf}
\caption{PJ04 R2 observed KEYGEN, SIGN, and VERIFY latency effects with validated uncertainty summaries.}
\label{fig:pj04_r2_latency_effects}
\end{figure*}

\begin{figure*}[htbp]
\centering
\includegraphics[width=0.88\textwidth]{pj04_r2_assets/PJ04_R2_FIG_SIGN_BLOCK_LATENCY.pdf}
\caption{PJ04 R2 block-level signing latency observations.}
\label{fig:pj04_r2_sign_blocks}
\end{figure*}

\begin{figure*}[htbp]
\centering
\includegraphics[width=0.88\textwidth]{pj04_r2_assets/PJ04_R2_FIG_VERIFY_BLOCK_LATENCY.pdf}
\caption{PJ04 R2 block-level verification latency observations.}
\label{fig:pj04_r2_verify_blocks}
\end{figure*}
% END PJ04_R2_INDIVIDUAL_OPT

\label{sec:individual}


% BEGIN CANONICAL PATHB POSTPATCH INTEGRATION V1
\subsection{Canonical post-patch Path-B WOTS+C evaluation}
\label{subsec:pathb-canonical-postpatch}

A fresh matched reference-core campaign was conducted after
canonical-verifier integration to establish the current
performance authority for the experimental Path-B WOTS+C
construction. The Standard XMSSMT control produced a
$4963$-byte signature, whereas canonical Path-B produced a
$4779$-byte signature, yielding an absolute reduction of
$184$ bytes ($3.707435\%$).

At the fresh-key-block level ($n=10$ paired blocks per
condition), Path-B reduced key-generation latency by
$4.841972\%$ relative to the Standard control
(pointwise paired-block bootstrap 95\% CI:
$-4.957407\%$ to $-4.741268\%$). Signing latency was reduced by
$4.750131\%$ (95\% CI:
$-4.858689\%$ to $-4.640895\%$).

Verification moved in the opposite direction. Canonical Path-B
increased verification latency by $5.190717\%$
(95\% CI: $+0.608855\%$ to $+10.058769\%$), from a mean of
$2.677961$ ms for the Standard control to $2.816966$ ms for
Path-B.

\begin{table}[htbp]
\centering
\caption{Canonical post-patch XMSSMT Path-B WOTS+C comparison
on Raspberry Pi 5. Confidence intervals are pointwise
operation-specific paired-block bootstrap intervals.}
\label{tab:pathb-canonical-postpatch}
\begin{tabular}{lrrrr}
\hline
Metric & Standard & Path-B & Effect & Interpretation \\
\hline
Signature size (B)
& 4963 & 4779 & $-184$ ($-3.707435\%$) & Path-B smaller \\
KEYGEN (ms)
& 2468.040906 & 2348.539062 & $-4.841972\%$ & Path-B lower \\
SIGN (ms)
& 4927.317295 & 4693.263271 & $-4.750131\%$ & Path-B lower \\
VERIFY (ms)
& 2.677961 & 2.816966 & $+5.190717\%$ & Path-B higher \\
\hline
\end{tabular}
\end{table}

These results establish an operation-specific
communication--latency trade-off rather than universal
performance superiority. The historical pre-patch
verification-speedup interpretation is superseded by this
canonical post-patch campaign.

Path-B remains an experimental, non-standard construction.
Security equivalence to conventional WOTS+/XMSSMT has not been
established, and standards-compliance claims are not authorized.
Consequently, Path-B is retained outside the
security-preserving Pareto pool.
% END CANONICAL PATHB POSTPATCH INTEGRATION V1

\subsection{Prior layer-adaptive XMSSMT evidence}
\label{subsec:prior_layer_adaptive_review}

The earlier layer-adaptive study supplies a twenty-configuration
dataset combining five ordered layer-height allocations with four
ordered WOTS parameter placements. This dataset is reused here as
prior evidence for the optimization framework and is not presented
as a new Paper~3 measurement campaign.

The frozen dataset demonstrates that configuration selection depends
on the selected objectives. Its strict Pareto summary contained three
configurations when minimizing mean signing latency, the signing
95th percentile, signature size, and serialized secret-key size.
Adding mean key-generation latency increased the strict frontier
to ten configurations; adding verification latency increased it
to twelve. These are results for the recorded configurations,
measurement windows, and objective definitions, rather than evidence
for a universally optimal profile.

The dataset combines original measurements with randomized repeated
measurements for the H14--6 configurations. Differences in observation
counts and measurement windows must therefore remain explicit.
This earlier-study frontier must also be distinguished from any
Paper~3 frontier involving checkpointing, streaming, hash backends,
or combined strategies.

\subsection{Layer-aware delta checkpointing}
\label{subsec:delta_checkpoint_review}

The XMSSMT source audit established that upper-layer WOTS signatures
were already cached in the baseline implementation. Consequently, adding
the same cache would duplicate existing functionality. The evaluated
extension instead used layer-aware dirty-state characterization and an
authenticated delta journal to reduce persistent checkpoint traffic.
For XMSSMT-SHA2\_20/2\_256, complete checkpoints were retained initially
and at lower-subtree rollover, while intermediate records encoded
contiguous changed byte ranges. The reconstructed standard secret-key
state and subsequent signatures remained byte-identical in the recorded
validation tests.

The state characterization covered three independent keys and 3078
transitions. Ordinary transitions had a median of 201 changed bytes,
whereas the three rollover transitions changed between 4482 and
4490 bytes. These measurements describe differences in byte values,
rather than an instrumented trace of memory stores.

The persistence evaluation used ten process sessions with alternating
variant order. Recorded mean ordinary commit latency was
16.380~ms for full-state persistence and 11.947~ms for delta persistence.
The archived analysis reported a mean latency reduction of 27.048\%
and an application-level write reduction of 91.911\%.
Recovery after 64 transitions was slower for the delta journal:
the recorded means were 259.898~microseconds for delta recovery and
49.115~microseconds for baseline recovery.

All fifteen recorded process-crash fault tests passed, including
subsequent signature-equivalence checks. These tests do not establish
power-loss durability. Rollback rejection assumes an external monotonic
floor, and application-level bytes must not be interpreted as physical
flash writes. The checkpoint optimization changes neither the standard
signature format nor the signature size.

\subsection{Compact state and persistence costs}
\label{subsec:compact_state_review}

The compact-state evaluation comprised ten process sessions for
XMSS-SHA2\_10\_256 and XMSSMT-SHA2\_20/2\_256. Mean snapshot sizes were
1118.058 and 4414.108 bytes, respectively, compared with raw serialized
states of 1377 and 6002 bytes. The archived session-level analysis
reported mean size savings of 18.805\% for XMSS and 26.456\% for XMSSMT.

Serialization overhead and durable persistence were evaluated
separately. Recorded mean pack/unpack costs were
5.819/4.138~microseconds for XMSS and
15.588/7.231~microseconds for XMSSMT.
Mean durable signing/export latency was 34.568~ms for XMSS,
compared with 6.701~ms for signing alone, and
37.430~ms for XMSSMT, compared with 9.183~ms for signing alone.
The archived analysis reported mean session-level overheads of
415.993\% and 307.812\%, respectively. Thus, reduced serialized-state
size did not imply a corresponding reduction in durable signing latency.

The provenance record reports successful correctness, boundary,
lower-tree rollover, and POSIX crash-durability validation.
However, the monotonic-floor backend was a file test double;
rollback-resistant hardware was not implemented. Snapshot confidentiality
was not provided, warm recovery was not a cold-reboot measurement,
and distributed multi-writer operation was outside the supported scope.

\subsection{Alternative hash backend}
\input{LH05_LH08_RPI5_HASH_RESULTS.tex}

\label{subsec:hash_backend}
% Earlier O4 results require separate integration.

\subsection{Streaming WOTS+ public-key reduction}
\label{subsec:o5_streaming_review}

The streaming implementation supplied WOTS+ chain elements directly to
the L-tree reduction instead of materializing the full WOTS+ public key.
Both variants used the same baseline OpenSSL SHA-256 backend. The
evaluation distinguished memory observations from latency measurements.

The memory pilot used a fixed pthread-stack write watermark rather than
total process resident memory. For XMSS-SHA2\_10\_256, the signing
watermark decreased from 10,928 to 9,488 bytes. For
XMSSMT-SHA2\_20/2\_256, it decreased from 11,824 to 10,384 bytes.
Both observations correspond to a reduction of 1,440 bytes. Verification
showed a reduction of 1,392 bytes for both profiles. These observations
describe the tested stack watermark and do not establish an equivalent
reduction in total process memory.

Latency was evaluated in ten process sessions, using session-level
Student-$t$ confidence intervals without outlier removal. Mean signing
latency changes were $+2.467\%$ for XMSS, with a 95\% interval of
$[-1.302,6.236]\%$, and $+0.941\%$ for XMSSMT, with an interval of
$[-0.426,2.309]\%$. All reported latency-change intervals included zero.
The results therefore support the observed memory benefit but do not
resolve a latency improvement or establish latency equivalence.

\subsection{Multicore execution}
\label{subsec:o6_multicore_review}

The multicore campaign compared serial execution with two-thread and
four-thread variants in ten process sessions. Each session included
three key-generation observations, twenty sequential signing
observations, and twenty verification observations per profile and
variant. Variant order was randomized using deterministic session
seeds. CPU affinity was CPU~0 for serial execution, CPUs~0--1 for two
threads, and CPUs~0--3 for four threads.

Four-thread signing achieved a mean session-level speedup of 1.953
for XMSS, with a 95\% interval of $[1.922,1.985]$, and 1.949 for
XMSSMT, with an interval of $[1.914,1.983]$. Four-thread verification
speedups were 1.655 and 1.664, respectively. Signing efficiency was
approximately 48.8\% for XMSS and 48.7\% for XMSSMT, indicating that
the measured acceleration did not scale linearly with thread count.

Signature-index allocation, private-key state mutation, and BDS state
mutation remained serialized. These results consequently describe
parallel work within the evaluated implementation and do not authorize
concurrent whole-signature operations on a shared mutable private key.

\subsection{Signature-size optimization}
\label{subsec:signature_size}
\subsubsection{Independent lossless signature encoding}
The O7 feasibility evaluation used 300 standard XMSS signatures
of 2500 bytes and 300 standard XMSSMT signatures of 4963 bytes.
ZLIB9, BZ2\_9, and LZMA\_XZ9 were evaluated using whole-signature
and component-wise encoding. Decoding required no previous
signature, cross-signature dictionary, or external decoder state.
All 3600 encoding observations reported byte-exact reconstruction.

The prespecified feasibility criterion required a mean size
reduction of at least 5\% and smaller encodings for at least 95\%
of signatures. No evaluated codec and mode met this criterion.
The best result for both schemes was whole-signature ZLIB9:
2521 bytes for XMSS and 4984 bytes for XMSSMT, representing
expansions of 0.84\% and 0.423\%, respectively. No tested encoding
was smaller than its original signature. The independent codec
path was therefore closed with a negative feasibility result.
This finding applies to the evaluated codecs and corpus; it does
not establish impossibility for every lossless encoding method.

The separate Windows WOTS+C construction comparison is reported
in the cross-platform results section.

\subsection{Cross-construction compression extension}
\label{subsec:cross_construction}
The O8 evaluation compared SLH-DSA-SHAKE-128s with the existing,
attributed SPHINCS+C-SHAKE-128s-simple construction. Their
signature sizes were 7856 and 6304 bytes, respectively: a
1552-byte difference, or 19.7556\%. Both used 32-byte public
keys and 64-byte secret keys. SPHINCS+C was evaluated as a
separate construction and was not classified as FIPS 205
SLH-DSA or as a newly proposed compression method.

The Raspberry Pi 5 campaign comprised ten paired sessions,
with one key generation, ten signatures, and ten verifications
per construction per session. Each construction ran in a fresh
process. Execution order alternated between sessions, giving
five sessions with each construction first. The protocol
specified CPU 0 affinity and the performance governor.
No observations were removed. Inference used paired session
means and a two-sided Student-$t$ interval on natural-log
latency ratios, with nine degrees of freedom.

Mean key-generation latency was 145.185 ms for SLH-DSA and
89.914 ms for SPHINCS+C. Mean signing latency was 1116.542
and 1335.806 ms, respectively; verification latency was
1.102 and 7.772 ms. The geometric SLH-DSA/SPHINCS+C ratios
and 95\% confidence intervals were 1.614694
[1.612707, 1.616683] for key generation, 0.837428
[0.799918, 0.876698] for signing, and 0.141809
[0.140383, 0.143250] for verification.

Thus, the smaller SPHINCS+C signatures accompanied faster
key generation and slower signing and verification in this
campaign. These differences describe the complete constructions
and parameterizations. They cannot be attributed solely to
compression, and they do not establish security equivalence
or a size reduction within standard SLH-DSA.

\input{MTLH05_MTLH08_XMSSMT_HASH_RESULTS.tex}

\input{SLH05_SLH08_DERIVED_HASH_RESULTS.tex}

\section{Combined Optimization}
\label{sec:combined}

\subsection{Direct evaluation of combined optimization}
\label{subsec:o9_combined_review}

The combined campaign evaluated two profiles and five implementation
arms in ten paired sessions. The primary comparison was C07\_BATCH4
against REF\_FIX01. Signing and verification used arithmetic means of
ten observations within each session; key generation used one
observation per session. Speedup was defined as reference latency
divided by candidate latency. The analysis used geometric means of
paired session ratios and 95\% bootstrap intervals with 10,000
replicates and seed 20260928.

For XMSS, primary speedups were 3.860 for key generation
(95\% interval $[3.816,3.903]$), 4.087 for signing
($[3.916,4.279]$), and 2.948 for verification
($[2.826,3.050]$). For XMSSMT, the corresponding estimates were
3.843 ($[3.814,3.866]$), 3.949 ($[3.676,4.207]$), and
2.676 ($[2.442,2.895]$).

These estimates came from direct combined-implementation measurements.
They were not obtained by multiplying speedups from separate module
campaigns. Across the full analysis, five comparison classifications
differed between mean-based primary and median-based sensitivity
analyses. Those disagreements must be reported alongside the
corresponding comparisons. The timing results alone do not establish
signature-size, serialized-state, memory, energy, or physical-storage
benefits.

\section{Experimental M3 XMSSMT Profile}
\label{sec:m3}

% This section is generated from the frozen M3 analysis package.

\input{m3_rpi5/M3_RPI5_RESULTS.tex}

\input{m3_rpi5/TABLE_M3_RPI5_PUBLICATION_RESULTS.tex}

\input{m3_rpi5/M3_RPI5_FIGURE_INSERTIONS.tex}

\input{m3_rpi5/M3_RPI5_DISCUSSION.tex}

\input{m3_rpi5/M3_RPI5_LIMITATIONS.tex}

\input{m3_rpi5/M3_RPI5_REPRODUCIBILITY_STATEMENT.tex}


\section{Joint Multi-Objective Analysis}
\label{sec:joint}

The evaluated strategies affect different deployment costs.
Signature bytes, transmitted packet bytes, serialized private state,
checkpoint traffic, working memory, and operation latency are distinct
quantities. A reduction in one must not be reported as a reduction in
another. Compact snapshots and delta checkpoints reduce state-related
costs without shortening the standard signature. Streaming reduces
working-buffer requirements, while parallel execution primarily targets
latency. Alternative hash backends change the experimental construction
and its measured execution cost.

The signature-size findings also require separate interpretation.
Independent lossless encoding increased the size of every tested
standard signature. WOTS+C produced smaller experimental XMSS and
XMSSMT signatures. SPHINCS+C provided a smaller signature in a
cross-construction comparison, with higher signing and verification
latencies on the Raspberry Pi 5. These results do not demonstrate
security-preserving compression of an unchanged standardized signature.

The earlier layer-adaptive study supplies a Pareto analysis for its
twenty recorded configurations. That frontier remains specific to
its configurations, measurement windows, and selected objectives.
The module campaigns in this manuscript use different protocols,
profiles, and observation counts. They therefore support a structured
comparison of trade-offs, but are not pooled into a new numerical
Pareto frontier. Implementation complexity is discussed qualitatively
where no common quantitative measure was collected.

Direct combined-implementation measurements provide evidence for
the combinations actually executed. Their effects cannot be inferred
by multiplying speedups from separate experiments. Deployment
selection should first specify the required construction and state
reliability model, then compare feasible implementations using the
costs relevant to the application. The present evidence does not
identify a universally best configuration.

\section{Cross-Platform Analysis}
\input{LH05_LH08_WINDOWS_HASH_RESULTS.tex}

\input{O10_WINDOWS_FASTCORE_RESULTS.tex}


% BEGIN PJ04_R2_CROSS_PLATFORM
The corrected PJ04 R2 evidence reported in this manuscript was obtained on the
Raspberry Pi~5 platform. These measurements therefore characterize the tested
Raspberry Pi~5 implementation and configuration and should not be interpreted
as demonstrating equivalent behavior on the Windows implementation or on
other platforms.
% END PJ04_R2_CROSS_PLATFORM

\label{sec:crossplatform}

% TODO:
% Integrate Windows results only after validation/freeze.
% Prefer normalized ratios where methodologies differ.

\section{Security and Correctness Considerations}
\label{sec:security}

The evaluation distinguishes functional correctness, implementation
robustness, and cryptographic security. Successful signing, verification,
message recovery, and rejection tests provide evidence for the exercised
implementation paths. They do not constitute a security proof.

The repaired Windows XMSS WOTS+C implementation was tested over its full
1024-signature lifetime. The tests checked private-state and signature indices,
message recovery, rejection after exhaustion, and 15 negative cases. A
short-input verification defect was corrected before these tests passed.
Buffer guards detected no overwrite into the guarded regions during the
tested executions; this does not establish general memory safety.

The Windows XMSSMT WOTS+C implementation generated and verified 1026
consecutive signatures, including the first height-10 subtree rollover.
Twenty-eight negative cases passed. Global exhaustion of the height-20
profile was not exercised. These results establish the tested rollover
behavior, rather than full-lifetime validation of that profile.

Stateful signing requires protection against reuse of a one-time signing
index. Sequential index progression in these tests does not establish
durable protection against process crashes, rollback, concurrent access,
or restoration of an old private-state copy. Claims about persistence and
recovery must therefore refer to the separate state-management experiments
and their stated failure model.

The alternative-hash and WOTS+C configurations are experimental
constructions. Retaining a 32-byte digest, preserving address encodings,
or producing a smaller signature does not by itself establish the security
level of the standard construction. Accordingly, the reported size and
latency results are implementation measurements, and security equivalence
with standardized XMSS or XMSSMT is not claimed. Evidence from these
stateful configurations also does not establish properties of SLH-DSA.

\section{Discussion}

% BEGIN CANONICAL PATHB DISCUSSION V1

The canonical Path-B result refines the interpretation of the
earlier WOTS+C measurements. Signature-size reduction and the
key-generation and signing advantages remain measurable in the
fresh matched campaign, but the canonical verification rule
introduces a verification penalty of $5.190717\%$. Therefore,
the final evidence supports a communication/performance
trade-off rather than a claim of universal speedup. This
distinction is methodologically important because the
canonical-verifier change affects the current verification
authority; pre-patch verification timings are not reused for
the final Path-B claim.

The Path-B evidence is also kept separate from the
security-preserving implementation-optimization frontier.
Although the construction is functionally validated under the
frozen experimental protocol, security equivalence and
standards compliance remain unestablished.

% END CANONICAL PATHB DISCUSSION V1


% BEGIN PJ04_R2_DISCUSSION
The PJ04 WOTS+C result illustrates a measured size--performance trade-off in
which the deterministic signature and FULL-packet footprints were each reduced
by 92~bytes while the observed mean KEYGEN, SIGN, and VERIFY latencies were
also lower for the evaluated PJ04\_WOTSC configuration under the measured
Raspberry Pi~5 conditions. This observation is specific to the tested
implementation, configuration, platform, and execution conditions; it does not
establish a causal performance effect or universal performance superiority.
% END PJ04_R2_DISCUSSION

\label{sec:discussion}

% TODO:
% Joint interpretation across O1--O13.

\section{Limitations}

% BEGIN PJ04_R2_LIMITATIONS
For PJ04 R2, CONTROL was measured before PJ04\_WOTSC, so temporal-order
effects cannot be excluded. In addition, the 100 signing and 100 verification
measurements per mode were nested within 10 execution blocks and must not be
treated as 100 independent experimental replicates. The validated bootstrap
independently resampled blocks within each mode and therefore does not imply
paired-block correspondence. Performance findings are specific to the measured
Raspberry Pi~5 implementation, configuration, and experimental conditions.

Thermal observations are included only as descriptive supporting platform
evidence and are not used to infer causal thermal effects.

\input{pj04_r2_tables/PJ04_R2_TABLE_PLATFORM_THERMAL.tex}

\begin{figure*}[htbp]
\centering
\includegraphics[width=0.88\textwidth]{pj04_r2_assets/PJ04_R2_FIG_THERMAL_SUMMARY.pdf}
\caption{PJ04 R2 descriptive thermal summary for the measured Raspberry Pi~5 runs.}
\label{fig:pj04_r2_thermal}
\end{figure*}
% END PJ04_R2_LIMITATIONS

\label{sec:limitations}

% TODO:
% Platform-specificity,
% experimental profiles,
% statistical replication,
% implementation-specific results,
% security-proof boundaries.

\section{Reproducibility}

% BEGIN PJ04_R2_REPRODUCIBILITY
The finalized PJ04 R2 results, figures, tables, provenance records, and
postvalidation evidence were assembled into a hash-verified publication
package and subsequently frozen in the validated publication-package evidence
chain. Manuscript integration uses those frozen artifacts without benchmark
re-execution, statistical re-execution, bootstrap re-execution, figure
regeneration, or numerical result recomputation.
% END PJ04_R2_REPRODUCIBILITY

\label{sec:reproducibility}

% TODO:
% Final Objective-1 artifact manifests,
% dataset hashes,
% source identities,
% analysis freezes.

\section{Conclusion}

% BEGIN CANONICAL PATHB CONCLUSION V1

For the canonical post-patch Path-B WOTS+C configuration,
XMSSMT signature size decreased from $4963$ B to $4779$ B
($184$ B; $3.707435\%$). Under the frozen Raspberry Pi 5
reference-core campaign, key-generation and signing latency
were lower by $4.841972\%$ and $4.750131\%$, respectively,
whereas verification latency was higher by $5.190717\%$.
Accordingly, this configuration is reported as an experimental
communication/performance trade-off and not as a universally
superior, security-equivalent, or standards-compliant
construction.

% END CANONICAL PATHB CONCLUSION V1


The study separates signature-size reduction from improvements in
state storage, checkpoint traffic, working memory, and execution
latency. The evaluated independent lossless codecs did not reduce
standard XMSS or XMSSMT signatures. Experimental WOTS+C configurations
reduced the evaluated XMSS signature from 2500 to 2408 bytes and the
XMSSMT signature from 4963 to 4779 bytes. These deterministic size
differences accompany implementation-specific timing results and
functional validation, rather than a demonstrated security-equivalence
result.

Compact state, delta checkpointing, streaming, parallel execution,
and directly measured combinations address complementary deployment
costs. Their usefulness depends on the selected profile, platform,
failure model, and application objective. The alternative-hash results
remain experimental, and the SPHINCS+C evaluation remains a comparison
between different constructions.

For Objective 1, the evidence establishes measured trade-offs and
tested implementation improvements. It does not establish that every
size-reducing experimental construction preserves the security
guarantees of its standardized counterpart. Functional correctness,
state reliability under the exercised failure model, and cryptographic
security must remain separately justified.

\label{sec:conclusion}



\backmatter

\bmhead{Acknowledgements}
% TODO.

\section*{Declarations}

\bmhead{Funding}
% TODO.

\bmhead{Competing interests}
The authors declare no competing interests.

\bmhead{Data availability}
% TODO: final repository/archive statement after Objective-1 freeze.

\bmhead{Code availability}
% TODO.

\bmhead{Author contributions}
% TODO.

\bibliography{references}

\end{document}
